V tejto práci sa zameriavame na problém odhadu intenzity slnečného žiarenia 
z fotografií oblohy. Pri návrhu modelov založených na konvolučných neurónových sieťach zohráva významnú úlohu nielen architektúra siete a jej hyperparametre,
 ale aj vlastnosti vstupných dát. Jedným z kľúčových faktorov je rozmer vstupného obrázka, ktorý ovplyvňuje 
 presnosť modelu a aj výpočtovú náročnosť trénovania.

Cieľom tejto práce je analyzovať vplyv rozmerov vstupných obrázkov na kvalitu regresného modelu a nájsť 
optimálne nastavenie pre rôzne architektúry konvolučných neurónových sietí. V rámci experimentov sú systematicky 
testované viaceré konfigurácie modelov a hyperparametrov, pričom výsledky sú vyhodnocované pomocou štandardných
 chybových metrík a časovej náročnosti trénovania.

Výsledkom práce je identifikácia optimálnych kombinácií architektúry a veľkosti vstupných obrázkov, ktoré poskytujú 
najlepší kompromis medzi presnosťou odhadu a výpočtovou efektívnosťou, ktorá zohrávala nemalú úlohu pri návrhu celej implementácie
vzhľadom na veľkosť datasetu.